\documentclass[a4paper]{article}

% Basic packages
\usepackage[fontset=fandol]{ctex}
\usepackage{amsmath, amssymb}
\usepackage{graphicx}
\usepackage[margin=2.5cm]{geometry}
\usepackage[most]{tcolorbox}
\usepackage{etoolbox}
\usepackage{listings}
\usepackage{hyperref}
\usepackage{booktabs}
\usepackage{subcaption}
\usepackage{float}
\usepackage{tikz}
\IfFileExists{pgfplots.sty}{
    \usepackage{pgfplots}
    \pgfplotsset{compat=1.18}
}{}

% Highlight boxes
\newtcolorbox{knowledgebox}[1]{
    enhanced, colback=blue!5!white, colframe=blue!75!black, colbacktitle=blue!75!black,
    coltitle=white, fonttitle=\bfseries, title=#1, attach boxed title to top left={yshift=-2mm, xshift=2mm},
    boxrule=1pt, sharp corners
}

\newtcolorbox{importantbox}[1]{
    enhanced, colback=yellow!10!white, colframe=yellow!80!black, colbacktitle=yellow!80!black,
    coltitle=black, fonttitle=\bfseries, title=#1, sharp corners
}

\newtcolorbox{warningbox}[1]{
    enhanced, colback=red!5!white, colframe=red!75!black, colbacktitle=red!75!black,
    coltitle=white, fonttitle=\bfseries, title=#1, sharp corners
}

\newtcolorbox{dialoguebox}[1]{
    enhanced, breakable, colback=green!4!white, colframe=green!45!black, colbacktitle=green!45!black,
    coltitle=white, fonttitle=\bfseries, title=#1, sharp corners
}

% Listings
\lstset{
    language=Python,
    basicstyle=\ttfamily\small,
    keywordstyle=\color{blue},
    stringstyle=\color{red!60!black},
    commentstyle=\color{green!60!black},
    breaklines=true,
    frame=single,
    numbers=left,
    numberstyle=\tiny\color{gray},
    captionpos=b,
    extendedchars=false
}

% Front-page metadata
\newcommand{\notetitle}{课程笔记：强化学习系统}
\newcommand{\noteauthors}{Bowen-mantou（整理）}
\newcommand{\notedate}{2026 年 10 月 23 日}
\newcommand{\videochannel}{常青藤名校课（B 站）}
\newcommand{\videopublishdate}{2026-09-09}
\newcommand{\videoduration}{66:27}
\newcommand{\videourl}{https://www.bilibili.com/video/BV1a6Yx62EH4?p=26}
\newcommand{\videocoverpath}{cover.jpg}

\begin{document}

\begin{titlepage}
\centering
\vspace{1.2cm}
{\huge\bfseries \notetitle\par}
\vspace{0.8cm}
{\large \noteauthors\par}
\vspace{0.3cm}
{\large CMU 11-768: AI Agents（2026 年秋季）\par}
{\large 主讲：CMU 11-768 课程组\par}
\vspace{0.3cm}
{\large \notedate\par}
\vspace{1.2cm}

\ifdefempty{\videocoverpath}{
{\small 视频封面因版权未收录，见下方视频链接。\par}
}{
\par
}

\vfill
\begin{tcolorbox}[width=0.9\textwidth, colback=black!2!white, colframe=black!60, sharp corners]
\textbf{视频作者/频道}：\videochannel\par
\textbf{发布日期}：\videopublishdate\par
\textbf{视频时长}：\videoduration\par
\textbf{视频链接}：
\ifdefempty{\videourl}{
未填写
}{
\href{\videourl}{\nolinkurl{\videourl}}
}
\par\textbf{PDF制作Skill}：\href{https://github.com/wdkns/wdkns-skills}{\nolinkurl{https://github.com/wdkns/wdkns-skills}}
\end{tcolorbox}
\end{titlepage}

\tableofcontents
\newpage

%% --- 正文内容开始 --- %%

\section{本讲定位：从算法走到「系统」}

\subsection{第三讲的任务：把算法装进机器}

三讲强化学习部分的最后一讲，主题从「算法」转向「系统」。前两讲讨论的是策略梯度、信用分配、异步训练与稳定性；这一讲要回答的是一个更工程化的问题：\textbf{如果让你从零搭一个能训练 LLM agent 的 RL 框架，你需要哪些部件？它们为什么这样切分？}

讲者给本讲的目标定得很具体：不是教你从头实现一个框架，而是\textbf{带你走一遍「如果你用 RL 框架训 agent、或者要 fork 修改某个框架时会遇到的东西」}。所以本讲的内容结构是「总览 → agentic RL 特有的坑 → 框架设计选择」，这也正好对应后面三节。

\subsection{一个 LLM RL 框架至少要支持哪些操作}

讲者先抛出一个问题：\textbf{要训练这样一个模型，框架里必须有哪些核心操作？}课堂给出的答案可以归成三类（图~\ref{fig:pg-primitives}）：

\begin{enumerate}
\item \textbf{前向（forward）}：既用来算 loss，也用来做 rollout 采样（agent 生成动作、与环境交互都要前向）；
\item \textbf{反向（backward）}：用来自动微分、计算梯度；
\item \textbf{优化器步（optimizer step）}：把算出来的更新真正写回模型权重。
\end{enumerate}

如果把策略梯度写成一句话，它正好把这三件事串起来：

$$\nabla_\theta J(\theta) = \mathbb{E}\Big[\sum_{t}\big(R(\tau)-b(h_t)\big)\,\nabla_\theta \log \pi_\theta(a_t \mid h_t)\Big]$$

逐项解释：$\theta$ 是策略参数；$\tau$ 是轨迹；$R(\tau)$ 是轨迹奖励；$b(h_t)$ 是基线（上一讲讨论过组内均值、价值函数等都是它的实现）；$\nabla_\theta \log \pi_\theta(a_t\mid h_t)$ 是「让这个动作更可能出现的梯度方向」。括号里的 $R(\tau)-b(h_t)$ 是权重，$\nabla_\theta\log\pi_\theta$ 这一项正是要反向传播的东西。

但对 \textbf{agentic} RL 来说，这三件事远远不够。讲者补充了三个 agent 特有的部件：

\begin{itemize}
\item \textbf{环境（environment）}：策略在环境里执行动作、拿回观测，所以框架必须能「step 环境」；
\item \textbf{奖励函数与优势（reward \& advantage）}：需要对采样到的轨迹做验证、打分、算优势；
\item \textbf{推理的结构（agent harness）}：agent 如何组织 prompt、工具调用、多轮交互，这决定了推理请求长什么样。
\end{itemize}

\subsection{为什么不能「用 PyTorch 手写」了事}

\begin{sloppypar}
很多人第一次接触这些框架时都会问：这些操作映射到 PyTorch 不就是 \texttt{model.forward}、\texttt{loss.backward}、\texttt{optimizer.step} 吗？为什么需要 SkyRL、Miles、veRL、PrimeRL、Slime 这么多框架？
\end{sloppypar}

讲者的回答是：\textbf{因为推理和训练是两种负载特性完全不同的 workload}。如果只写一个朴素的 PyTorch 循环，你会在两边同时损失性能：

\begin{itemize}
\item \textbf{推理侧}：自回归地一个 token 一个 token 解码，是 \textbf{decode 为主}的负载；不需要 autograd、不需要保存计算图，但需要大量显存给 \textbf{KV cache}；
\item \textbf{训练侧}：整条序列都已知，可以\textbf{并行}处理（更像 prefill）；必须保存激活值和计算图才能反向传播。
\end{itemize}

两种负载的最优实现方式差别太大，于是几乎所有成熟框架都采用\textbf{双引擎}架构：一个专门做推理（例如 vLLM、SGLang），一个专门做训练（例如 Megatron、FSDP2），中间通过「权重同步」和「轨迹数据」两个通道连接（图~\ref{fig:two-workloads}）。

\subsection{一个训练步的完整生命周期}

把双引擎拼起来之后，一个训练步长这样（图~\ref{fig:training-step}）：

\begin{enumerate}
\item \textbf{sample}：推理引擎与 agent harness 交互，在环境里执行动作、拿到观测，采出一条条轨迹；
\item \textbf{verify}：环境或验证器给出奖励信号（测试是否通过、答案是否正确等）；
\item \textbf{forward}：训练引擎把整条轨迹并行地重新前向一遍，\textbf{保留激活与计算图}；
\item \textbf{backward}：反向传播算梯度；
\item \textbf{optimizer step + weight sync}：更新参数，并把新权重同步回推理引擎，开始下一轮。
\end{enumerate}

讲者特别提醒：推理引擎和训练引擎是\textbf{同一个模型的两套实现}，它们的 kernel 或前向计算方式可能略有不同，因此即使权重完全相同，采样出来的分布也会有细微差异——也就是说，框架在设计上\textbf{天然带一点 off-policy}。所以上一讲讨论的重要性采样、裁剪等 off-policy 校正在这里不是可选项，而是必需品。

\subsection{本章小结}

\begin{itemize}
\item LLM RL 框架的核心操作是前向、反向、优化器步；agentic RL 还要加上环境交互、奖励验证、harness 结构。
\item 不能用单一 PyTorch 循环替代框架：推理与训练的负载特性（decode vs prefill、是否保存计算图、显存去向）根本不同，所以采用推理/训练双引擎。
\item 一个训练步是「采样 → 验证 → 前向 → 反向 → 更新与权重同步」的循环；两套引擎的计算差异天然带来轻微 off-policy，需要显式校正。
\end{itemize}

\section{双引擎架构：推理与训练各自的优化}

\subsection{推理侧：KV cache 是显存的主要去向}

自回归解码时，每生成一个新 token 都要对全部历史做注意力；如果每步都重算 prefix 的 K/V，计算量会随长度平方增长。\textbf{KV caching} 就是把这些算过的 key/value 存进显存，用「显存」换「计算」。

讲者给了一个数量级估算，用来说明推理侧显存压力有多大：\textbf{Llama-3-8B、BF16 精度下，64 条 32K token 的轨迹，KV cache 大约要 256 GB}（讲者口头估算，具体数值随实现而变，细节待核对）。而且这还只是一个 8B 的小模型；在 RL 场景里，batch size 可以非常大，KV 数据会更多。

这也解释了为什么双引擎架构有优势：如果训练和推理共用一个引擎，你既想把激活值留在显存里做反向，又想给 KV cache 留足空间，显存会严重冲突；分开之后，推理引擎可以几乎把显存都花在 KV cache 上。

\subsection{推理侧：跨请求的 prefix 复用甚至更重要}

KV cache 解决的是「一个请求内部」的重复计算；agentic RL 里还有大量「\textbf{请求之间}」的重复：一次 agent 轨迹的每一步都是一个独立的 LLM 请求，它们共享同一段系统 prompt、同一个任务 prompt，甚至前一步的完整历史就是后一步的 prefix。

\begin{importantbox}{一个关键事实}
「轨迹的每一步是一个独立的 LLM 请求，不是一次连续的推理。」——所以在 agentic RL 里，如果引擎不支持 \textbf{prefix caching}（前缀缓存），同一段 prefix 会被反复 prefill 很多次，浪费是成倍的。
\end{importantbox}

主流的实现方式有两种：SGLang 的 \textbf{RadixAttention}（用基数树管理前缀、复用粒度细）和 vLLM 的 \textbf{automatic prefix caching}（按 block 内容哈希复用）。讲者把「跨请求复用」列为推理引擎对 agentic RL 最重要的能力之一。

\subsection{推理侧：continuous batching、chunked prefill 与 PD 分离}

\textbf{continuous batching}（连续批处理）解决的是「长度不齐带来的空转」（图~\ref{fig:cb}）：朴素批处理要等整批里最长的请求做完，短请求完成后槽位就空转；连续批处理则在\textbf{每个 decode step 之后重新做调度决策}——刚完成的请求立刻退出 batch，腾出的容量马上给等待中的新请求。

但实际的推理负载不只有 decode，还有 \textbf{prefill}（把 prompt 一次性并行处理）。长 prompt 的 prefill 会阻塞很多 decode 请求，所以现代引擎会做两件事：

\begin{itemize}
\item \textbf{chunked prefill}：把长 prompt 切成块，与 decode 交错处理，避免一次性长时间占住 GPU；
\item \textbf{disaggregated prefill/decode（PD 分离）}：把一部分 GPU 专门用于 prefill、另一部分专门用于 decode。讲者解释它的另一个好处是\textbf{异构硬件适配}：例如 NVIDIA 服务器里既有适合高吞吐 prefill 的硬件（讲者口述提到 Groq LPU、Blackwell/Rubin 一类，名称听记）、也有适合 decode 的 GPU，分离后可以为不同阶段挑选更合适的硬件；代价是 prefill 与 decode 之间要传递 KV（通常通过队列/服务完成）。
\end{itemize}

\subsection{推理侧：投机解码在 RL 里的特殊问题}

\textbf{speculative decoding}（投机解码，图~\ref{fig:spec}）用一个更小更快的 draft 模型先连续猜几个 token，再让大模型并行验证；配合拒绝采样，可以在\textbf{保持与原大模型分布等价}的前提下加速。

讲者指出它在 RL 训练中有一个特殊陷阱：训练初期 acceptance rate 可能很高，但 RL 在不断改模型，\textbf{基座模型一漂移，draft 模型的接受率就会掉下来}。常见的应对是用辅助目标（SFT/蒸馏）在训练过程中\textbf{持续更新 draft 模型}，让它跟上基座。

\subsection{训练侧：激活检查点与训练侧的 prefix 复用}

训练侧的核心矛盾是\textbf{激活值显存}。反向传播需要前向时算过的激活，全存下来显存会爆；\textbf{activation checkpointing}（激活检查点，图~\ref{fig:ckpt}）的做法是只保留每层输入、丢弃中间激活，反向时再\textbf{选择性重算}缺失的激活。

讲者给了量化代价：这相当于多做一次前向，\textbf{大约多花 33\% 的时间}，但省下大量显存，是「显存受限时的好朋友」。实践中还分 \textbf{full} 与 \textbf{selective}（只对部分层做检查点）两种模式。

推理侧的 prefix 复用在训练侧同样存在，而且更复杂：训练要保留激活并做反向。讲者提到的代表方案是 Snowflake 的 \textbf{Arctic RL} 里的 \textbf{Zorro}，以及其中使用的 \textbf{dynamic tree attention}：把一个 batch 的公共前缀组织成\textbf{前缀树}，沿树做 DFS——先对公共祖先做前向，进入某一分支做前向与反向，反向完成后立刻\textbf{弹出计算图、释放激活}，再回退走另一分支。这样显存中同时保留的激活\textbf{永远不超过一条根到叶的路径}，既复用了 prefix，又没有把显存撑爆。

\subsection{训练侧：五种并行分片策略}

模型大到单卡放不下时，训练引擎要把模型/数据切到多卡上。讲者按「切什么」把主流方法过了一遍，这里汇总成表~\ref{tab:shard}。

\begin{table}[H]
\centering
\small
\begin{tabular}{@{}p{0.13\textwidth}p{0.24\textwidth}p{0.26\textwidth}p{0.25\textwidth}@{}}
\toprule
\textbf{策略} & \textbf{切分对象} & \textbf{通信特点} & \textbf{适用场景} \\
\midrule
DP（数据并行） & 全局 batch 切成 micro-batch，每卡一份完整模型 & 每步梯度 all-reduce，一次、可跨机 & 默认首选；有「按 token 预算装箱」的 micro-batch 划分算法 \\
PP（流水线并行） & 按层切成 stage，不同 rank 放不同层 & 相邻 stage 传激活；用 micro-batch 填充流水 & 层数多、单卡装不下整模型 \\
EP（专家并行） & MoE 的不同专家放到不同 rank & 路由带来的 all-to-all & MoE 模型 \\
TP（张量并行） & 把单层计算（矩阵乘/注意力头）切开 & 通信量大，通常只在节点内 NVLink 上 & 单层都放不下时；应尽量少用 \\
CP（上下文并行） & 按序列维切，各 rank 持一段 token & 注意力需要交换 KV（环形传递） & 超长上下文 \\
\bottomrule
\end{tabular}
\caption{五种并行分片策略对照。讲者的经验法则是：\textbf{尽量把并行度花在 DP 上，减少 TP 之类通信密集的方案}。}
\label{tab:shard}
\end{table}

讲者给的经验法则值得单独记住：\textbf{优先最大化数据并行，减少张量并行的比例}——因为 TP 的通信开销与层内依赖最强，通常限制在单机内；而 DP 可以跨机扩展。分片策略对推理同样重要，只是推理侧的约束从「激活」变成了「KV cache」。

\subsection{引擎交互之一：权重同步为什么能只传 20--30 GB}

两个引擎之间的第一个通道是\textbf{权重同步}。直觉上这件事很贵：一个 1T 参数模型，BF16 下约 \textbf{2 TB}；如果每次更新都把全量权重广播到所有推理 rank，根本不可行（讲者也提到这正是「跨集群训练 + 存到 S3 再下载」这类做法的现实动机）。

但讲者指出一个关键观察：\textbf{RL 的更新非常稀疏}。在 BF16 精度下，两次更新之间\textbf{约 99\% 的参数是逐比特相同的}——因为每条轨迹的监督信号很少、更新幅度极小，很多参数在 BF16 下根本没有变化。于是可以\textbf{只同步 delta}：需要传输的数据从 2 TB 降到大约 \textbf{20--30 GB}，就完全可行了。

\begin{knowledgebox}{为什么「99\% 参数没变」是合理的}
BF16 的尾数只有 7 位（讲者在 Q\&A 里确认），表示精度有限；RL 里每条轨迹的信号又很稀疏，很多参数的更新量小于 BF16 能表示的最小变化，于是它们在数值上保持不变。讲者对这个现象的实用判断是：\textbf{如果某次更新小到被 BF16 丢掉、但它重要，下一次训练步会出现更大的更新补回来}；而用 32 位训练在显存上又不可行——所以 BF16 训练 + delta 同步是现实的最优解。
\end{knowledgebox}

同步完权重还有一层工程细节：训练引擎的分片布局与推理引擎的分片布局通常不同，所以需要\textbf{重新分片（re-shard）}，让权重符合推理引擎的切分方式。

\subsection{引擎交互之二：MoE 的「路由重放」（R3）}

第二个不显眼但很要命的交互是 \textbf{MoE 路由}。MoE 前向时要为每个 token 选择专家；由于实现细节差异、甚至只是路由打分出现平局，\textbf{训练引擎选中的专家可能与推理引擎当时选中的不同}——这就让训练的前向与推理的前向不一致，产生额外的 off-policy。

讲者给的解法非常直接（图~\ref{fig:r3}）：在 rollout 阶段\textbf{记录每个 token 的路由决策}，随轨迹数据一起传给训练引擎；训练引擎前向时\textbf{丢弃自己的路由结果、强制重放这些记录}。这个技术被称为 \textbf{Rollout Routing Replay（R3）}，讲者说它对 MoE 训练稳定性帮助很大，并给了一句很实用的排查建议：\textbf{如果你在训 MoE 模型时发现 loss 或 reward 不涨，先检查有没有开 R3}。

\subsection{本章小结}

\begin{itemize}
\item 推理侧优化围绕 KV cache 展开：KV caching、prefix caching（RadixAttention / automatic prefix caching）、continuous batching、chunked prefill、PD 分离、投机解码（RL 下要在线更新 draft）。
\item 训练侧优化围绕激活显存展开：activation checkpointing（约 33\% 时间换显存）、训练侧 prefix 复用（Arctic RL/Zorro 的 dynamic tree attention）、五种并行分片。
\item 引擎交互有两件事：权重同步（利用更新稀疏性，2 TB $\rightarrow$ 20--30 GB 的 delta）与 MoE 路由重放（R3，MoE 训练稳定性必备）。
\end{itemize}

\section{Agentic RL 的「尖锐细节」（The Sharp Bits）}

\subsection{序列扩展：为什么「压缩上下文」的 harness 更难训}

agent harness 有两种典型设计（图~\ref{fig:seqext}）：

\begin{itemize}
\item \textbf{harness A（追加式）}：每一步都把新的动作/观测\textbf{追加}到历史后面，历史前缀永远保持不变；
\item \textbf{harness B（压缩式）}：每做几步就把之前的动作与观测\textbf{总结压缩}成一段 summary，再从 summary 继续。生产系统里很常见（讲者点名 Claude Code、OpenHands 都做 context compaction）。
\end{itemize}

课堂投票里多数人认为 B 更好（因为上下文更短），但讲者的答案是：\textbf{从训练效率看 A 更好}——只有在序列实在太长、必须压缩时才用 B。

原因在 batch 的组织方式（图~\ref{fig:seqbatch}）：harness A 的前缀完全连续，\textbf{整条轨迹可以当作一个数据点}进 batch，一次 prefill 就够；harness B 一旦压缩，就把轨迹断成了两段独立序列，\textbf{必须拆成两个数据点}——于是任务 prompt 被重复 prefill 两次，计算浪费。要缓解这一点，就得依靠前面讲过的 training 侧 prefix 复用技巧（如 dynamic tree attention 把多个数据点组织成树来复用计算）。

\subsection{Chat template 的坑：推理链被悄悄删掉}

比上下文压缩更隐蔽的问题是 \textbf{chat template}。讲者以 Qwen3 / Qwen3.5 的默认 chat template 为例（图~\ref{fig:template}）：开启 reasoning 后，模型会先输出一段思维链再给动作；但当新的 user turn 或 harness 输出进来时，\textbf{默认模板会把历史里的 reasoning token 删掉}。

这在工程上可以理解（节省上下文窗口、避免干扰），但在 RL 训练里它等价于「每一步都做一次上下文压缩」，\textbf{直接破坏了序列扩展性质}：prefix 不连续、历史被改写。于是他给出了一个很实用的结论：\textbf{很多开源 RL 框架会自带定制的 chat template（现在常被称为 renderer）}，用来保证前缀不丢；一些训练服务还提供开关（例如「跨轮保留 reasoning」「完全不要 reasoning」）。讲者还提到课程作业里训练小 RL 框架时会提供保留前缀的 Qwen3.5 模板（课程内部资料）。

\subsection{Token-in, token-out：不要用字符串在引擎之间传数据}

早期框架在推理引擎与训练引擎之间传轨迹时，很自然地用\textbf{字符串}（因为 OpenAI 风格的 chat completions API 就是字符串消息，prompt 工程也都在字符串上做）。讲者说，\textbf{几乎所有框架后来都发现这是错的}，正确的做法是 \textbf{token-in, token-out}：保存并传递\textbf{精确的 token 序列}（图~\ref{fig:tito}）。

原因有两层：

\begin{itemize}
\item \textbf{字符串侧的微小变换会放大成 tokenization 差异}：训练引擎的 chat template 只要插入一个空格或换行，重新 tokenize 出来的 token 序列就可能与推理时不同，于是训练和采样用的根本不是同一串 token——直接 off-policy。
\item \textbf{tokenization 不是一一对应的}：文本到 token 是确定的，但 token 到文本再回 token 不唯一。讲者举了一个例子：某个词在推理时被模型输出成「1 个 token + 5 个 token」的组合，但重新 tokenize 同一段字符串可能得到「4 号 token」——同一个字符串对应不同 token 序列，训练目标就被悄悄改掉了。
\end{itemize}

讲者说这类问题在早期框架里造成过很多莫名其妙的 \textbf{loss 尖刺与训练不稳定}，现在几乎所有 RL 框架都会保存精确 token 并直接传给训练引擎。

\subsection{异步 RL 在 agent 训练里为什么更重要}

上一讲从「稳定性」角度讲了异步 RL；这一讲从「系统」角度再讲一遍（图~\ref{fig:asyncrl}）：agent 轨迹长度差异极大，同步训练时 GPU 大量时间在等最长的 rollout（straggler）；异步训练让推理与权重更新\textbf{重叠}，更长的轨迹可以拖到下一批处理，短轨迹则可以先被用上更新。

讲者还给了两个具体的系统设计事实：

\begin{itemize}
\item \textbf{同步训练的常见做法是 co-locate}：8 张卡先全部用来训练，再让训练引擎「睡眠」、把同一批卡切给推理引擎，交替使用；\item \textbf{异步训练可以拆分 GPU 配比}：因为推理与训练同时活跃，可以给更吃紧的一侧更多卡。讲者说他在实践中常用 \textbf{3:1} 的比例（例如 6 个节点做推理、2 个节点做训练），因为长程 agent 任务里 rollout 往往是瓶颈。
\end{itemize}

\subsection{In-flight 权重更新：生成到一半，权重换了}

异步 RL 有一个很「反直觉」的操作：推理引擎正在为某些请求生成时，训练完成了新一轮更新，需要\textbf{中途换权重}。讲者把这类操作称为 \textbf{in-flight weight updates}，并指出推理引擎为此专门演进出了接口。

以 SGLang 的文档为例（图~\ref{fig:pause}），引擎提供 \textbf{pause generation} 接口，其中 \textbf{retract} 模式很有用：它不只是暂停生成，还会\textbf{清空 KV cache}——因为参数更新后，prefix 的 KV 已经\textbf{过期}，必须用新模型重算。讲者总结：vLLM 与 SGLang 都在为 RL 框架的需求持续演进。

\subsection{复杂 harness：递归多 agent 与「一棵树怎么进 batch」}

agentic RL 的最后一个系统难点来自 harness 本身越来越复杂。讲者展示了他近期论文里训练\textbf{递归多 agent 系统}的推理模式：一个 root agent 把任务委派给子 agent，子 agent 还可以继续委派——这来自「递归语言模型」的思路。再叠加前面说的上下文压缩，推理结构就变成一棵树。

于是框架必须回答一个没有标准答案的问题：\textbf{整棵树的一次 rollout 算一个数据点，还是把每个子 agent 的轨迹当成独立数据点进 batch？}不同的选择会影响 prefix 复用、loss mask、优势广播与 batch 组织。讲者没有给唯一答案，只是把它作为「agent harness 复杂化带来的系统设计问题」摆出来。

\subsection{本章小结}

\begin{itemize}
\item 序列扩展：追加式 harness（A）对训练更友好，压缩式（B）会把一条轨迹拆成多个数据点、产生重复 prefill；chat template 若删除历史 reasoning，等价于偷偷做压缩。
\item 引擎之间必须 \textbf{token-in, token-out}：字符串重 tokenize 会导致 off-policy 与 loss 尖刺。
\item 异步 RL 是 agent 训练的系统默认选择：overlap、拆分 GPU 配比（常见 3:1）、in-flight 权重更新（pause/retract 并清 KV cache）。
\item 递归多 agent 与 context compaction 让「如何组织 batch」成为新的开放设计问题。
\end{itemize}

\section{框架的系统设计选择}

\subsection{Harness-agnostic：用一个「中间人」隔离 harness 与训练框架}

agent harness 生态极其多样（讲者点名 Codex、Claude Code、OpenHands、PydanticAI（Py）、Hermes、Prime Agent 等，且每天都在增加）。如果每加一个 harness 就要改 RL 框架，维护成本会爆炸。

讲者介绍的流行做法是 \textbf{harness-agnostic}：在 harness 与推理引擎之间放一个\textbf{中间人/代理服务}，它\textbf{模拟 harness 已经在用的推理 API}（chat completions 或 responses API），把请求转发给推理引擎、把响应转发回 harness；同时\textbf{在中间把训练需要的一切记录下来}（logprob、R3 的路由数据等）。这样 harness 侧零改动，框架只维护一个代理；这个代理甚至可以打包成一个「推理服务」对外暴露 endpoint。

\subsection{SPMD：一个程序跑遍所有卡，简单但不灵活}

讲者接着讨论分布式代码的组织方式。最直观的是 \textbf{SPMD}（Single Program Multiple Data）：所有 GPU rank 跑\textbf{同一份 Python 程序}，靠 \texttt{rank} 做条件分支，让不同 rank 执行不同角色（伪代码见下方）。

代码长这样（示意，非原文）：

\begin{lstlisting}[caption={SPMD 风格的伪代码（据讲者幻灯片整理）：所有 rank 跑同一程序，按 rank 分组执行训练或推理分支。}]
train = dist.new_group([0, 1, 2, 3])
infer = dist.new_group([4, 5, 6, 7])

if rank < 4:
    batch = next_batch_shared(train)
    forward_backward(train)
    optimizer_step()
else:
    while True:
        refresh_weights_if_ready(infer)
        trace = rollout(prompt, infer)
        enqueue_trajectory(trace, infer)
\end{lstlisting}

讲者对这个模式的评价很中肯：优点是\textbf{直观、全局可见、改起来直接}（完整逻辑在一个程序里）；缺点是在大规模、多角色、异构硬件的场景里\textbf{非常不灵活}。他举的实例是 AReaL：它在 1.0 之前采用 SPMD，之后转向了下面要讲的服务化架构。

\subsection{Everything-as-a-service：控制器 + 标准协议}

另一种更现代的模式是「\textbf{一切皆服务（everything as a service）+ 单控制器}」：训练服务、推理引擎、权重同步服务、agent harness、环境服务各自作为独立服务存在，一个\textbf{轻量的 CPU 控制器}负责编排（图~\ref{fig:svc}）。

讲者列出的优点很具体：

\begin{itemize}
\item \textbf{分离关注点}：每个组件都极其复杂（推理、训练各有自己的优化体系），拆开后可以独立构建、维护、升级；
\item \textbf{标准协议}：harness 与工具/环境之间可以用 \textbf{MCP}，推理侧用 chat completions / responses API；
\item \textbf{轻量控制器}：编排循环可以是纯 CPU 的，不需要 GPU、不引入复杂依赖；
\item \textbf{模块化可替换}：只要协议一致，推理引擎可以换 vLLM/SGLang，训练框架可以换，甚至 harness 也可以换；
\item \textbf{弹性与异构}：可以同时挂多个推理服务（甚至 AMD 与 NVIDIA 混用），需要时动态加机器。
\end{itemize}

讲者举了两个例子。一是 \textbf{AstroFlow}（Haizhong Zheng 等，2026）——它把 rollout 做成服务：推理服务 + agent + 环境都可以在不同 GPU 上、按需加入；权重同步服务把权重推给所有推理服务。二是 \textbf{Tinker API}：一个标准的训练 API，最初由 Thinking Machines 为其托管训练服务提出，现在 Miles、SkyRL 等开源框架也实现它——\textbf{同一份训练代码可以在托管服务与自建集群之间迁移}。

\subsection{服务化带来的三种新能力}

讲者进一步说明，服务化不只是「解耦好看」，它打开了三类以前很难做的事：

\begin{enumerate}
\item \textbf{多策略训练（multi-policy）}：不同 rollout/agent 可以跑不同权重甚至不同架构的模型，因为可以为不同模型分配不同训练器与推理引擎。这对多 agent 异构训练（不同 agent 用不同模型）很关键。
\item \textbf{在线 RL（训练后部署）}：把线上部署的 harness 也当成数据来源，定期把真实使用数据送进训练服务、更新模型再上线。讲者举的例子是 Cursor：其 tab 模型与 Composer 模型都有持续收集用户数据、更新并部署新版本的流程；据其博客，使用 auto 功能时大约\textbf{每 5 小时}就可能用到新的 checkpoint（讲者口述引用博客，细节待核对）。
\item \textbf{多租户 / 多 LoRA 训练}：托管训练服务面对很多用户、每人一套 LoRA。Tinker 最早做出这种模式：一个训练 client 对应一次 run，一个采样 client 对应某个 LoRA 权重检查点；底层用 \textbf{Grouped-GEMM / SGMV} 等专用 kernel，把多个 LoRA 的推理与反向\textbf{并发}执行。
\end{enumerate}

讲者给了这组对比的结论：8 个并发 LoRA run 相比 8 个串行 run，\textbf{整体吞吐明显更高}——单个 run 因为共享算力会更慢，但设备利用率上去了。他还提到开源实现（SkyRL 及 Trajectory Labs 一侧的实现）可以自建同样的能力，而 Tinker 是把这套「多租户 + 解耦训练/采样 API」做出名的服务。

\subsection{本章小结}

\begin{itemize}
\item harness-agnostic：在 harness 与推理引擎之间加一个模拟标准推理 API 的代理，顺带采集训练数据；harness 零改动。
\item SPMD 简单直接但不够灵活（AReaL 1.0 前用过）；服务化 + 轻量控制器带来分离关注点、标准协议、可替换、弹性与异构硬件、多策略、在线 RL、多租户多 LoRA 等能力。
\item 落地建议：先对齐标准 API（Tinker API、chat completions），再按「推理服务 / 训练服务 / 权重同步 / 环境 / harness」拆服务；控制器保持轻量。
\end{itemize}

\section{总结与延伸}

\subsection{讲者的收尾：一整讲其实在讲三件事}

讲者在结尾没有给「标准答案框架」，而是把整讲收敛成三件事：

\begin{enumerate}
\item \textbf{两种负载必须分开优化}：推理（decode-heavy、KV cache）与训练（prefill-style、激活与反向）各有自己的优化体系；双引擎是主流答案。
\item \textbf{引擎交互是新的故障源}：权重同步（利用稀疏性传 delta）、MoE 路由重放（R3）、token-in/token-out、off-policy 校正——这些细节做错，训练就会出现莫名其妙的尖刺与不涨点。
\item \textbf{框架设计在向「服务化 + 标准协议」迁移}：harness-agnostic、轻量控制器、可替换组件、多租户多 LoRA、在线 RL，都是这个方向上的能力。
\end{enumerate}

\subsection{一张表把系统问题映射到排查手段}

\begin{table}[H]
\centering
\small
\begin{tabular}{@{}p{0.24\textwidth}p{0.32\textwidth}p{0.34\textwidth}@{}}
\toprule
\textbf{症状} & \textbf{可能的系统原因} & \textbf{排查/修复} \\
\midrule
训练 loss 或 reward 出现尖刺、不稳定 & 引擎间传字符串导致重 tokenize 不一致；MoE 路由不一致 & 改为 token-in/token-out；检查 R3 是否开启 \\
推理/训练整体吞吐低、GPU 空转 & 同步训练等 straggler；batch 未连续调度；KV cache 分配不合理 & 上异步 RL；continuous batching；prefix caching \\
权重同步太慢 & 每次广播全量权重 & 只传 delta（BF16 下约 99\% 参数不变），或走对象存储 \\
训练显存爆 & 激活未检查点；prefix 重复计算 & activation checkpointing（约 33\% 时间代价）；training 侧 prefix 复用（dynamic tree attention） \\
上下文越长越慢、训练效率下降 & harness 压缩历史 / chat template 删 reasoning / 轨迹被拆成多个数据点 & 改用追加式 harness；定制保留前缀的 renderer；序列级 batch 组织 \\
框架难扩展、加一个 harness 要改代码 & SPMD 单程序耦合 & 服务化 + 标准协议；harness-agnostic 代理 \\
\bottomrule
\end{tabular}
\caption{本讲系统问题的症状—原因—修复对照表。}
\end{table}

\subsection{课堂最后一个问答：BF16 更新被丢掉，会不会损失信号？}

有人问：如果 BF16 精度下很多参数更新小到「没有变化」，是不是在丢训练信号？讲者的回答包含三点：

\begin{itemize}
\item 每条轨迹的监督信号本来就很少（outcome reward 更极端），所以单次更新幅度极小，很多参数在 BF16 下确实不变；
\item 如果这个更新真的重要，\textbf{下一次训练步会出现更大的更新}把它补回来；
\item 用 32 位训练在显存上不可行，所以「BF16 训练 + delta 同步」是现实的取舍，而不是疏忽。
\end{itemize}

\begin{importantbox}{本讲最值得带走的一条工程原则}
agentic RL 的系统问题，绝大多数不是「某个组件不够快」，而是\textbf{两个组件之间的语义没有对齐}：权重怎么同步、路由怎么重放、token 怎么传、prefix 怎么复用、更新到什么程度算「同一个策略」。把接口和语义对齐，性能问题往往就变成可诊断的问题。
\end{importantbox}

\subsection{延伸阅读}

\begin{itemize}
\item 推理引擎：vLLM 的 PagedAttention 与 automatic prefix caching；SGLang 的 RadixAttention 与 pause generation 接口（本讲 in-flight 权重更新的例子来源）。
\item 训练系统：Megatron-LM 的 TP/PP/EP 并行；FSDP2 的 ZeRO 式分片；Arctic RL（Snowflake）与 Zorro 的 dynamic tree attention。
\item 框架与 API：Tinker（Thinking Machines）的多租户多 LoRA 训练 API；AReaL 的 SPMD $\rightarrow$ 服务化演进；AstroFlow（Haizhong Zheng 等，2026）的服务化 rollout 设计。
\end{itemize}

\end{document}
